\section{Integración causal y validación numérica}
\label{bre:num:seccion}

La integración determina si una semilla obtenida mediante balance armónico o
criterios geométricos alcanza un régimen persistente del sistema estudiado.
Para interpretar esa trayectoria deben fijarse el campo, los parámetros, el
orden, la condición inicial, el terminal inferior y el tratamiento de la
memoria. La comparación entre métodos conserva estos datos; los cambios de
paso, horizonte o aproximación del núcleo se evalúan por separado.

\subsection{Discretización de Volterra mediante ABM--PECE}
\label{bre:num:pece}

Considérese el problema conmensurable
\begin{equation}
 {}^{\mathrm C}D_{t_0}^{q}X(t)=F(t,X(t)),\qquad X(t_0)=X_0,
 \quad 0<q\leq1,
 \label{bre:num:pvi}
\end{equation}
con $F$ continua y localmente Lipschitz respecto del estado. En un intervalo
de existencia, su forma integral es
\begin{equation}
 X(t)=X_0+\frac1{\Gamma(q)}
 \int_{t_0}^{t}(t-s)^{q-1}F(s,X(s))\,\dd s.
 \label{bre:num:volterra}
\end{equation}
El esquema de Adams--Bashforth--Moulton con secuencia PECE integra este núcleo
exactamente y aproxima el campo dentro de cada intervalo
\cite{DiethelmFordFreed2002}. En la malla $t_n=t_0+nh$, sea
$F_j=F(t_j,X_j)$. Sustituir el campo por $F_j$ en $[t_j,t_{j+1}]$ da
\[
 \int_{t_j}^{t_{j+1}}(t_{n+1}-s)^{q-1}\,\dd s
 =\frac{h^q}{q}\bigl[(n+1-j)^q-(n-j)^q\bigr].
\]
Como $q\Gamma(q)=\Gamma(q+1)$, el predictor es
\begin{align}
 X_{n+1}^{\mathrm P}
 &=X_0+\frac{h^q}{\Gamma(q+1)}\sum_{j=0}^{n}b_{j,n+1}F_j,
 \label{bre:num:predictor}\\
 b_{j,n+1}&=(n+1-j)^q-(n-j)^q.\nonumber
\end{align}

Para el corrector se interpola linealmente:
\[
 F(t_j+hu,X(t_j+hu))\approx(1-u)F_j+uF_{j+1},
 \qquad 0\leq u\leq1.
\]
Con $r=n+1-j$, las contribuciones del intervalo son
$h^q[L_q(r)F_j+R_q(r)F_{j+1}]$, antes del factor $1/\Gamma(q)$. Las
integrales de las dos funciones base se calculan mediante
\begin{align}
 I_q(r)&=\int_0^1(r-u)^{q-1}\,\dd u
       =\frac{r^q-(r-1)^q}{q},\nonumber\\
 R_q(r)&=\int_0^1u(r-u)^{q-1}\,\dd u\nonumber\\
       &=rI_q(r)-\frac{r^{q+1}-(r-1)^{q+1}}{q+1},\nonumber\\
 L_q(r)&=I_q(r)-R_q(r).
 \label{bre:num:pesos-integrales}
\end{align}
La expresión de $R_q$ sigue del cambio $v=r-u$, que transforma $u$ en
$r-v$. Al reunir denominadores se obtiene
\begin{align*}
 q(q+1)L_q(r)&=(q+1-r)r^q+(r-1)^{q+1},\\
 q(q+1)R_q(r)&=r^{q+1}-(r+q)(r-1)^q.
\end{align*}
El nodo inicial recibe sólo $L_q(n+1)$; un nodo interior $j$ recibe
$R_q(n+2-j)+L_q(n+1-j)$; el extremo nuevo recibe $R_q(1)$. En un nodo
interior, los términos proporcionales a $r^q$ se combinan como
$[(q+1-r)-(r+1+q)]r^q=-2r^{q+1}$. Por tanto, los pesos normalizados son
\begin{align}
 a_{0,n+1}&=n^{q+1}-(n-q)(n+1)^q,\nonumber\\
 a_{j,n+1}&=(n-j+2)^{q+1}+(n-j)^{q+1}\nonumber\\
 &\hspace{1.1cm}-2(n-j+1)^{q+1},\quad 1\leq j\leq n.
 \label{bre:num:pesos-corrector}
\end{align}
Además, $q(q+1)R_q(1)=1$ y
$q(q+1)\Gamma(q)=\Gamma(q+2)$. Evaluar el campo desconocido del extremo en
el predictor produce
\begin{equation}
 \begin{split}
 X_{n+1}=X_0+\frac{h^q}{\Gamma(q+2)}\Bigl[
 &F(t_{n+1},X_{n+1}^{\mathrm P})\\
 &+\sum_{j=0}^{n}a_{j,n+1}F_j\Bigr].
 \end{split}
 \label{bre:num:corrector}
\end{equation}
La secuencia consiste en predecir, evaluar el campo, corregir y evaluar de
nuevo $F_{n+1}=F(t_{n+1},X_{n+1})$. Este último valor se incorpora a la
historia. Las sumas comienzan siempre en $j=0$: la implementación directa de
$N$ pasos requiere $O(N^2)$ operaciones y $O(N)$ almacenamiento para una
dimensión de estado fija.

\subsection{Límite entero y significado del predictor}
\label{bre:num:limite-entero}

Para $q=1$, $b_{j,n+1}=1$, $a_{0,n+1}=1$ y los pesos interiores del
corrector son $2$. Las fórmulas recuperan las cuadraturas compuestas
rectangular izquierda y trapezoidal de la ecuación integral; el extremo
derecho de la segunda se evalúa en el predictor. En particular,
\[
 X_{n+1}^{\mathrm P}=X_0+h\sum_{j=0}^{n}F_j
 =X_n^{\mathrm P}+hF_n\qquad(n\geq1).
\]
La recurrencia parte del predictor anterior, mientras que el estado corregido
satisface
\begin{align*}
 X_n&=X_0+\frac h2\Bigl[F_0+2\sum_{j=1}^{n-1}F_j
                         +F(t_n,X_n^{\mathrm P})\Bigr],\\
 X_n^{\mathrm P}-X_n&=\frac h2\bigl[F_0-F(t_n,X_n^{\mathrm P})\bigr].
\end{align*}
Por ello $X_{n+1}^{\mathrm P}$ no es, en general, el paso local
$X_n+hF_n$ iniciado en el estado corregido. Para $\dot x=x$, $x_0=1$, el
primer corrector da $x_1=1+h+h^2/2$, mientras que
\[
 x_2^{\mathrm P}-(x_1+hx_1)=-h^2/2.
\]
Este ejemplo distingue las identidades exactas de los pesos de una
identificación incorrecta con un método local de un paso.

\subsection{Continuación con historia y reinicio del sistema objetivo}
\label{bre:num:continuacion}

Sea $F_\eta$ una familia que une un sistema auxiliar, en $\eta=0$, con el
campo objetivo, en $\eta=1$. Se prescribe $\eta(t)=\eta_k$ en
$[t_k,t_{k+1})$ y se resuelve un único problema
${}^{\mathrm C}D_{t_0}^{q}X=F_{\eta(t)}(X)$. Cambiar $\eta$ modifica el
campo a partir de ese instante y conserva las evaluaciones anteriores.

La contribución de la historia previa al nivel $k$ es
\begin{equation}
 H_k(t)=\frac1{\Gamma(q)}\int_{t_0}^{t_k}
             (t-s)^{q-1}F_{\eta(s)}(X(s))\,\dd s.
 \label{bre:num:historia-heredada}
\end{equation}
Separar la integral en $t_k$ y usar $X(t_k)=X_0+H_k(t_k)$ da
\begin{equation}
 \begin{split}
 X(t)={}&X(t_k)+H_k(t)-H_k(t_k)\\
 &+\frac1{\Gamma(q)}\int_{t_k}^{t}
            (t-s)^{q-1}F_{\eta_k}(X(s))\,\dd s.
 \end{split}
 \label{bre:num:continuacion-integral}
\end{equation}
Un reinicio en $t_k$ elimina $H_k(t)-H_k(t_k)$. Para $q=1$, el núcleo es
constante y esa diferencia se anula. Para $0<q<1$, el peso de toda la
historia cambia al avanzar $t$, de modo que el estado terminal no basta para
reconstruir la evolución heredada.

Si el parámetro cambia por saltos, el campo dependiente del tiempo es
continuo por tramos. La formulación integral conserva sentido, pero una
estimación de interpolación basada en derivadas temporales acotadas no puede
atravesar esos saltos sin tratamiento adicional. Los tiempos \(t_k\) se
incluyen como fronteras de la partición de cuadratura. En cada lado se usa la
rama correspondiente; una interpolación lineal del campo debe distinguir
\(F_{\eta_{k-1}}(X(t_k))\) de \(F_{\eta_k}(X(t_k))\). La trayectoria es
continua, aunque los dos valores del campo puedan diferir. Si una
discretización usa un único valor nodal, el defecto causado por esa elección
forma parte del control numérico de la transición.

Al finalizar la continuación puede definirse un experimento autónomo nuevo,
\[
 {}^{\mathrm C}D_0^q Z(t)=F_1(Z(t)),\qquad Z(0)=X(t_{\mathrm{fin}}).
\]
Su objetivo es comparar esa condición inicial con los sondeos bajo un mismo
terminal inferior y una misma inicialización de memoria. La continuación
transporta una historia; el reinicio establece la trayectoria de referencia
del campo objetivo. Cambiar el orden $q$ también exige especificar un nuevo
problema, pues cambia el núcleo de la ecuación integral.

\subsection{Comparación con la recurrencia ADM local}
\label{bre:num:adm-local}

Wu et al. distinguen la recurrencia local \(X_{n+1}=\Phi_h(X_n)\) de un
predictor--corrector que utiliza la historia. Su ecuación (27) conserva
explícitamente el incremento de esa historia, y los autores advierten la
pérdida de memoria histórica del ADM por subintervalos en integraciones
largas \cite{Wu2023ArctanChua}. La diferencia se precisa restando la ecuación
de Volterra en \(t_n\) de la misma ecuación en \(t_{n+1}\):
\begin{multline}
 X(t_{n+1})-X(t_n)=\mathcal M_n\\
 +\frac1{\Gamma(q)}\int_{t_n}^{t_{n+1}}
 (t_{n+1}-s)^{q-1}F(s,X(s))\,\dd s,
 \label{bre:num:incremento-exacto}
\end{multline}
\begin{equation}
 \begin{split}
 \mathcal M_n&=\frac1{\Gamma(q)}\int_{t_0}^{t_n}B_n(s)F(s,X(s))\,\dd s,\\
 B_n(s)&=(t_{n+1}-s)^{q-1}-(t_n-s)^{q-1}.
 \end{split}
 \label{bre:num:incremento-memoria}
\end{equation}
La integral sobre $[t_n,t_{n+1}]$ describe el intervalo nuevo.
$\mathcal M_n$ repondera los intervalos anteriores y equivale a
\(H_n(t_{n+1})-H_n(t_n)\) en
\eqref{bre:num:continuacion-integral}, para un campo fijo. Se anula cuando
\(q=1\), pero en general no lo hace para \(0<q<1\). Una expansión ADM
local, aunque aumente su orden de truncamiento, no recupera ese término si
se reinicia usando sólo \(X_n\).

El campo constante \(F(t,x)=v\ne0\), con \(t_0=0\), proporciona un
contraejemplo explícito. La solución causal exacta y su incremento son
\begin{align*}
 X(t)&=X_0+\frac{t^q}{\Gamma(q+1)}v,\\
 X(t_{n+1})-X(t_n)
 &=\frac{h^q}{\Gamma(q+1)}[(n+1)^q-n^q]v.
\end{align*}
Resolver exactamente cada intervalo reiniciado produciría, en cambio,
\(\widehat X_{n+1}-\widehat X_n=h^qv/\Gamma(q+1)\). La contribución omitida
es
\[
 \mathcal M_n=\frac{h^q}{\Gamma(q+1)}
               [(n+1)^q-n^q-1]v.
\]
En un horizonte fijo \(T=Nh\), el reinicio repetido da
\[
 \widehat X_N=X_0+\frac{Th^{q-1}}{\Gamma(q+1)}v,
 \qquad
 X(T)=X_0+\frac{T^q}{\Gamma(q+1)}v.
\]
Para \(q<1\), la primera expresión ni siquiera converge a la segunda al
reducir \(h\). El error procede de omitir la historia, aun con solución
local exacta. Por ello una reconstrucción del algoritmo ADM publicado y una
integración causal con terminal fijo cumplen funciones de comparación
diferentes. La transferencia de una semilla entre ambas rutas requiere
reintegrarla bajo el operador de Caputo y el convenio de memoria elegidos,
y evaluar de nuevo su destino.

\subsection{Defecto integral, error de memoria y refinamiento}
\label{bre:num:error}

Una curva numérica continua $\widetilde X$ se contrasta con el problema de
Caputo mediante su defecto de Volterra,
\begin{equation}
 \begin{split}
 \rho_h(t)={}&\widetilde X(t)-X_0\\
 &-\frac1{\Gamma(q)}\int_{t_0}^{t}
       (t-s)^{q-1}F(s,\widetilde X(s))\,\dd s.
 \end{split}
 \label{bre:num:defecto}
\end{equation}
Supóngase que la solución y la aproximación permanecen en una región donde
$F$ es Lipschitz respecto del estado, con constante $L$. Restar sus
ecuaciones y tomar normas produce, para $t_0=0$,
\[
 u(t)\leq A+\frac L{\Gamma(q)}\int_0^t(t-s)^{q-1}u(s)\,\dd s,
 \quad u=\norm{X-\widetilde X},
\]
donde $A=\norm{\rho_h}_{\infty,[0,T]}$. La estimación se deriva iterando
la integral fraccionaria $I^q$. La identidad beta
\[
 \int_0^t(t-s)^{a-1}s^{b-1}\,\dd s
 =t^{a+b-1}\frac{\Gamma(a)\Gamma(b)}{\Gamma(a+b)}
\]
para $a,b>0$ implica $I^aI^b=I^{a+b}$ y
$I^{jq}1=t^{jq}/\Gamma(jq+1)$. Después de $N$ sustituciones,
\[
 u(t)\leq A\sum_{j=0}^{N-1}\frac{L^jt^{jq}}{\Gamma(jq+1)}
             +L^N(I^{Nq}u)(t).
\]
El último término está acotado por
$\norm{u}_\infty(LT^q)^N/\Gamma(Nq+1)$ y tiende a cero. Así se obtiene
la desigualdad de Gronwall fraccionaria \cite{Gomoyunov2018Convex},
\begin{equation}
 \sup_{0\leq t\leq T}\norm{X(t)-\widetilde X(t)}
 \leq A E_q(LT^q).
 \label{bre:num:gronwall}
\end{equation}
Para $q=1$, $E_1(z)=\exp z$. La cota es condicional a la permanencia en
la región indicada; una certificación requiere justificar también esa
permanencia y acotar el defecto, incluida su cuadratura.

Una aproximación de memoria debe controlar el núcleo en una norma adecuada.
Para $K_q(t)=t^{q-1}/\Gamma(q)$, una suma finita
$K_m(t)=\sum_{\ell=1}^m w_\ell\exp(-\eta_\ell t)$, con pesos y tasas
positivos, es acotada en cero y no aproxima uniformemente la singularidad si
$q<1$. Se separa $[0,\delta]$ y se exige
\begin{equation}
 \begin{split}
 \varepsilon_K(T)&=\int_0^T|K_q(s)-\widetilde K_m(s)|\,\dd s\\
 &\leq\varepsilon_{\mathrm{loc}}+(T-\delta)\varepsilon_\infty,
 \qquad T\geq\delta.
 \end{split}
 \label{bre:num:error-nucleo}
\end{equation}
Aquí $\widetilde K_m$ incluye el tratamiento local elegido,
$\varepsilon_{\mathrm{loc}}$ es su error integral en $[0,\delta]$ y
$\varepsilon_\infty$ acota el error en $[\delta,T]$. Una corrección local
exacta hace $\varepsilon_{\mathrm{loc}}=0$, pero conserva la necesidad de
almacenar la historia reciente.

Si $\norm{F}\leq M$ en la región y el defecto de cuadratura es a lo sumo
$\varepsilon_{\mathrm{quad}}$, sumar y restar $K_q*F(X_m)$ separa el error
de campo del error de núcleo. La segunda contribución está acotada por
$M\varepsilon_K$, por lo que \eqref{bre:num:gronwall} da
\begin{equation}
 \norm{X-X_m}_{\infty,[0,T]}
 \leq(M\varepsilon_K+\varepsilon_{\mathrm{quad}})E_q(LT^q).
 \label{bre:num:error-soe}
\end{equation}
Esta estimación tiene horizonte finito. Por ejemplo, sin corrección local,
las variables $z_\ell(t)=\int_0^t e^{-\eta_\ell(t-s)}F(s,X_m(s))\,\dd s$
satisfacen
\[
 \dot z_\ell=-\eta_\ell z_\ell+F(t,X_m),\qquad
 X_m=X_0+\sum_\ell w_\ell z_\ell.
\]
Para $F(x)=-ax$, $a>0$, el equilibrio proyectado es
$x_m^*=x_0/(1+a\sum_\ell w_\ell/\eta_\ell)$, pues
$z_\ell^*=-ax_m^*/\eta_\ell$. Puede diferir del límite cero del problema
de Caputo. Por ello la concordancia en $[0,T]$ no establece por sí sola
equivalencia de atractores o cuencas al llevar $T$ a infinito.

\paragraph{Regularidad y error de cuadratura.}
Los pesos exactos no eliminan la singularidad del arranque. Si \(F\) es
\(C^1\) cerca de \((t_0,X_0)\), la acotación del campo en
\eqref{bre:num:volterra} da \(X(t_0+s)-X_0=O(s^q)\). Por Lipschitz,
\(F(t_0+s,X(t_0+s))=F_0+O(s^q)\), donde \(F_0=F(t_0,X_0)\). Una nueva
sustitución en la integral produce
\begin{equation}
 X(t_0+s)=X_0+\frac{s^q}{\Gamma(q+1)}F_0+O(s^{2q}).
 \label{bre:num:arranque}
\end{equation}
Si \(0<q<1\) y \(F_0\ne0\), el cociente
\((X(t_0+s)-X_0)/s\) tiene norma no acotada cuando \(s\to0^+\). Así, la
suavidad del campo no implica una solución con derivadas temporales acotadas
en el terminal inferior.

Para precisar el papel de la regularidad, supóngase provisionalmente que
\(g(t)=F(t,X(t))\) tiene derivadas temporales acotadas
\(\norm{g'}\leq M_1\) y \(\norm{g''}\leq M_2\) en toda la ventana.
La interpolación constante tiene error a lo sumo \(M_1h\). La interpolación
lineal \(\ell_j\) en un intervalo de longitud \(h\) satisface
\[
 \norm{g(t)-\ell_j(t)}
 \leq\frac{M_2}{2}(t-t_j)(t_{j+1}-t)
 \leq\frac{M_2h^2}{8}.
\]
Esta cota se obtiene integrando dos veces el resto, nulo en ambos extremos;
el núcleo de esa representación tiene integral absoluta
\((t-t_j)(t_{j+1}-t)/2\). Multiplicar ambos errores por el núcleo positivo y
usar \(\int_0^T K_q(s)\,\dd s=T^q/\Gamma(q+1)\) da
\[
 \varepsilon_{\mathrm P}\leq\frac{M_1T^q}{\Gamma(q+1)}h,
 \qquad
 \varepsilon_{\mathrm L}\leq\frac{M_2T^q}{8\Gamma(q+1)}h^2.
\]
Estas son cotas de cuadratura sobre valores exactos de la trayectoria.
Sustituir el campo del extremo por su evaluación en el predictor añade a la
segunda cota, a lo sumo,
\[
 \frac{Lh^q}{\Gamma(q+2)}\varepsilon_{\mathrm P}.
\]
La consistencia resultante es \(O(h^2+h^{1+q})\) bajo esas hipótesis. Pasar
a un orden global del método requiere controlar además la propagación de
los errores nodales. Si las derivadas de \(g\) no son acotadas cerca de
\(t_0\), las cotas anteriores no son aplicables con constantes uniformes;
se necesita estimar el tramo inicial, introducir correcciones de arranque o
analizar una malla graduada. En las evaluaciones a paso uniforme se conserva
la evidencia de refinamiento sin atribuirle automáticamente el orden del
régimen suave.

\paragraph{Comparaciones a resolución creciente.}
El refinamiento temporal compara el mismo problema con $h,h/2,h/4$, usando
una malla común para medir el error de trayectoria en una ventana corta.
Después compara observables postransitorios al ampliar $T$ y el tiempo de
descarte. Si se aproxima la memoria, se varían además sus modos o longitud,
con error de núcleo controlado. La persistencia de una etiqueta clasifica
el destino al nivel de resolución evaluado; para estimar un orden de
convergencia se necesitan diferencias cuantitativas de la misma magnitud.
Con la distancia de \eqref{bre:num:distancia}, en una ventana corta fija \(I_c\) se define
\[
 E_h=\max_{t\in I_c}d_S(X_h(t),X_{h/2}(t)).
\]
Si \(E_h\simeq Ch^p\), con \(C>0\), entonces
\(E_h/E_{h/2}\simeq2^p\) y
\(p\simeq\log_2(E_h/E_{h/2})\). Esta estimación requiere diferencias no
nulas y un régimen asintótico de discretización; no se aplica a trayectorias
caóticas arbitrariamente largas. Perturbar $X_0$ a paso fijo evalúa otra propiedad: la robustez frente
a condiciones iniciales.

\subsection{Escalas, crecimiento tangente y comparación de nubes}
\label{bre:num:metricas}

La escala de las variables se fija mediante
$S=\operatorname{diag}(s_1,\ldots,s_n)$, $s_i>0$, y la distancia
\begin{equation}
 d_S(x,y)=\norm{S^{-1}(x-y)}_2.
 \label{bre:num:distancia}
\end{equation}
Para comparar escalas temporales se fija además \(\tau_*>0\). Con
\(t=t_0+\tau_*\theta\), \(Y(\theta)=S^{-1}X(t_0+\tau_*\theta)\), el
cambio de variable en Volterra da
\[
 {}^{\mathrm C}D_0^qY(\theta)
 =\tau_*^q S^{-1}F(t_0+\tau_*\theta,SY(\theta)).
\]
El factor \(\tau_*^q\) proviene del producto
\(\tau_*^{q-1}\tau_*\) entre el núcleo y el diferencial. Por tanto,
\(\widehat h=h/\tau_*\), \(\widehat T=T/\tau_*\) para una duración física \(T\); las escalas se
fijan antes de clasificar destinos. Si se usa una tasa espectral
\(\lambda\) de una linealización fraccionaria, su dimensión es
\(\mathrm{tiempo}^{-q}\), y una escala dimensionalmente coherente es
\(\tau_*=|\lambda|^{-1/q}\), para \(\lambda\ne0\).

Con $S=I$ se recupera la distancia euclidiana. En una variable angular se
usa la diferencia mínima módulo $2\pi$. Si se identifica un grupo finito
$G$ de simetrías, la expresión
$d_{S,G}(x,y)=\min_{g\in G}d_S(x,gy)$ define una métrica entre órbitas
cuando cada $g$ es una isometría. Para $g(x)=A_gx+b_g$, esto exige
$A_g^{\mathsf T}S^{-2}A_g=S^{-2}$. La invariancia permite cambiar
representantes, y la desigualdad triangular sigue de
$d_S(x,ghz)\leq d_S(x,gy)+d_S(gy,ghz)$.

En orden entero, para un campo \(C^1\), la matriz tangente satisface
\(\dot Y=DF(X(t))Y\), \(Y(0)=I\). En un campo continuo y \(C^1\) por
regiones, la integración variacional usa el Jacobiano de cada región y
resuelve los cruces de sus interfaces. En un cruce transversal de un campo
continuo, la matriz tangente es continua; los contactos tangenciales
requieren un análisis específico. Sobre una ventana de longitud $T$, el
crecimiento en la norma escalada se calcula con
$Y_S=S^{-1}YS$. Puesto que
$\norm{Y_Sv}_2^2=v^{\mathsf T}Y_S^{\mathsf T}Y_Sv$, los factores
principales de amplificación son los valores singulares $\sigma_i(Y_S)$.
Los exponentes de tiempo finito son
\begin{equation}
 \Lambda_i(T)=\frac1T\log\sigma_i(Y_S(T)).
 \label{bre:num:ftle}
\end{equation}
La acumulación de factores triangulares mediante reortogonalización QR
proporciona estimadores de los exponentes de Lyapunov
\cite{Benettin1980a}; sus valores finitos dependen de la ventana y de la
base inicial. Un exponente máximo positivo indica expansión durante el
intervalo observado. El horizonte heurístico de predicción
$T_{\mathrm{pred}}\simeq\lambda_1^{-1}
\log(\varepsilon_{\mathrm{adm}}/\varepsilon_0)$ se obtiene al resolver
$\varepsilon_0e^{\lambda_1T}=\varepsilon_{\mathrm{adm}}$, para
$0<\varepsilon_0<\varepsilon_{\mathrm{adm}}$ y $\lambda_1>0$
\cite{VallejoSanjuan2017Predictability}. Después de ese horizonte son
especialmente relevantes las distribuciones de retorno, los tiempos de
escape y los observables estadísticos convergentes.

En Caputo, la ecuación variacional conserva el mismo terminal inferior y
la historia del campo tangente. Para $F\in C^1$ en una región que contenga
las trayectorias próximas durante $[0,T]$, derivar Volterra respecto del
dato inicial produce
\begin{equation}
 Y(t)=I+\frac1{\Gamma(q)}\int_0^t(t-s)^{q-1}J(s)Y(s)\,\dd s.
 \label{bre:num:tangent-memory}
\end{equation}
Aquí $J(s)=DF(X(s))$. La derivación se justifica tomando cocientes
incrementales: el teorema
fundamental del cálculo expresa la diferencia de $F$ mediante la integral
de $DF$ sobre el segmento entre las soluciones. Su convergencia uniforme
en el compacto y la desigualdad integral de Gronwall hacen converger esos
cocientes a la única solución lineal anterior. En regiones no suaves se
necesita justificar separadamente esa diferenciabilidad.

Si en un tiempo $\tau$ se factoriza $Y(\tau)=QR$ con $R$ invertible,
el cambio de base $Z(t)=Y(t)R^{-1}$ satisface
\begin{equation}
 Z(t)=R^{-1}+\frac1{\Gamma(q)}
 \int_0^t(t-s)^{q-1}J(s)Y(s)R^{-1}\,\dd s.
 \label{bre:num:qr-memory}
\end{equation}
Por tanto, se transforman también el dato inicial tangente y todos los
valores históricos de $JY$. Conservar sólo $Z(\tau)=Q$ y reiniciar la
integral en $\tau$ elimina una contribución no constante del pasado.
Para $q=1$ esa contribución sí se reduce al estado tangente en $\tau$.
La identidad exige invertibilidad del factor utilizado; ésta debe
comprobarse, pues una matriz fundamental fraccionaria no hereda sin prueba
todas las propiedades de la de una EDO.
Al comparar algoritmos variacionales fraccionarios, incluidos los de
\cite{DancaKuznetsov2018Lyapunov}, se comprueba esta compatibilidad antes
de interpretar sus exponentes.

Para una nube objetivo $A=\{a_k\}$ y una cola de trayectoria
$B=\{b_j\}_{j=1}^{N}$, un criterio de contacto es
\begin{equation}
 D_{90}(B,A)=Q_{0.90}\bigl(\{\min_k d_S(b_j,a_k)\}_j\bigr)
 \leq\varepsilon_A.
 \label{bre:num:contacto}
\end{equation}
El percentil se calcula después del submuestreo declarado. Si las distancias
ordenadas son $d_{(1)}\leq\cdots\leq d_{(N)}$, se interpola entre las
posiciones $i+1$ e $i+2$, con $i=\lfloor0.90(N-1)\rfloor$ y peso
$0.90(N-1)-i$; para $N=1$ se toma la única distancia. Es un estadístico
dirigido: una sonda puede aproximar parte de la nube sin recorrerla toda.
La referencia, las escalas y el umbral se fijan antes del sondeo. La
concordancia geométrica se contrasta además con el destino terminal y la
estabilidad de los observables al variar paso y horizonte.

\subsection{Sondeos, salidas y censura temporal}
\label{bre:num:sondeos}

Alrededor de cada equilibrio $e$ se distinguen muestras volumétricas y
superficiales. En dimensión $n$, una muestra uniforme en una bola euclidiana
de radio $r$ es $x_0=e+rU^{1/n}d$, donde $U$ es uniforme en $[0,1]$ y
$d$ es independiente y uniforme en la esfera unidad. En efecto,
$\Pr(rU^{1/n}\leq a)=(a/r)^n$, la fracción de volumen contenida en el
radio $a$. En el sondeo superficial se usa $x_0=e+rd$; una distribución
determinista de direcciones proporciona un conteo a una resolución angular
fija. Los radios, direcciones y cantidades de muestras determinan el alcance
de la ausencia de contactos.

Para separar recurrencia y escape, sea $D$ una región cerrada de observación de un
flujo $\varphi_t$. El tiempo de primera salida es
\begin{equation}
 \tau_D(x_0)=\inf\{t\geq0:\varphi_t(x_0)\notin D\},
 \qquad \inf\varnothing=\infty.
 \label{bre:num:salida}
\end{equation}
Una integración hasta $T$ proporciona el dato censurado
\begin{equation}
 \widetilde\tau_D=\min(\tau_D,T),\qquad
 \Delta_D=\mathbf1_{\{\tau_D\leq T\}}.
 \label{bre:num:censura}
\end{equation}
En el cálculo se especifican la función de evento, su dirección de cruce,
la interpolación y la tolerancia. Si localmente $D=\{g\leq0\}$, una salida
transversal cumple $g(X(\tau_D))=0$ y
$\nabla g\cdot F>0$. Para $g_j\leq0<g_{j+1}$, el cero de la interpolación
lineal es
\[
 \widehat\tau_D=t_j-
 \frac{g_j}{g_{j+1}-g_j}(t_{j+1}-t_j).
\]
Si el cruce es único y la derivada temporal del evento es al menos
$\nu>0$, un error uniforme $\varepsilon_g$ de la función de evento implica
un error de raíz a lo sumo $\varepsilon_g/\nu$, más la tolerancia del
buscador. Esta estimación incluye el error de trayectoria y de
interpolación. Sin control entre muestras pueden omitirse salidas y
reentradas. La ausencia de escape detectado se considera censura derecha;
un fallo del integrador constituye un resultado distinto.

Para datos iniciales en un conjunto medible $U\subset D$ y una medida con
$0<\mu(U)<\infty$, la supervivencia se define por
\begin{equation}
 S_D(t)=\frac{\mu\{x_0\in U:\tau_D(x_0)>t\}}{\mu(U)}.
 \label{bre:num:supervivencia}
\end{equation}
Con $N$ semillas de igual peso y un horizonte común, para $t<T$ se estima
mediante $N^{-1}\sum_j\mathbf1_{\{\widetilde\tau_D^{(j)}>t\}}$, porque
$\min(\tau_D,T)>t$ equivale a $\tau_D>t$ en ese intervalo. Los intervalos
de confianza requieren un modelo de muestreo; los conteos de una malla
determinista no lo proporcionan por sí solos. Para una trayectoria de Caputo, \eqref{bre:num:salida} se aplica a su
evolución causal con el terminal inferior fijado, sin presuponer un
semiflujo sobre el estado proyectado. Se especifica la familia de historias
o reinicios: dos estados proyectados iguales pueden
tener escapes diferentes si sus memorias difieren. La persistencia, la
atracción y la separación respecto de los equilibrios requieren controles
propios; una trayectoria larga o un exponente positivo de tiempo finito no
sustituye esos argumentos.
